\documentclass[../main.tex]{subfiles}
\graphicspath{{\subfix{../tables_figures_sync}}{\subfix{../tables_figures_sync/}}}

\begin{document}

\section*{How to Use This Template}

This is a tutorial template for writing manuscripts in \LaTeX{} using Overleaf. Each section is meant to provide an example of that section of a paper and contains some pieces that demonstrate some features you may find useful for writing papers.

\subsection*{Template Structure}

This project is organized into multiple files:

\begin{itemize}
    \item \texttt{main.tex} -- The ``conductor'' that pulls everything together
    \item \texttt{doc\_setup.tex} -- Package imports and formatting settings
    \item \texttt{analysis\-values.tex} -- Number plugging commands and values
    \item \texttt{sections/} -- Individual section files (this modular approach keeps things much more organized!)
    \item \texttt{images/} -- Store all figures or images here
    \item \texttt{references.bib} -- Your bibliography file
\end{itemize}

\bigskip % one option for space between paragraphs

\subsection*{What you will learn in each section:}

\begin{itemize}
    \item \textbf{Abstract \& Introduction}: Basic text formatting, sections, citations
    \item \textbf{Methods}: Tables (simple and complex), equations
    \item \textbf{Results}: Figures -- both standalone and wrapped in text, using number plugging tex file (\texttt{analysis-values.tex})
    \item \textbf{Discussion}: Cross-referencing figures, tables, and sections
    \item \textbf{Appendix}: Supplementary materials
\end{itemize}

\vspace{1em} % another option that you can customize

\subsection*{Tip: comments}
This template includes custom commands for leaving comments. Try them out:

\lbw{This is Lauren's comment -- appears in blue.}

\jac{This is a comment from JAC -- appears in red.}

\mak{This is MAK's comment -- appears in magenta.}

These are defined in \texttt{doc\_setup.tex} and you can add more for other collaborators. You can also use the track changes/comment features, which we will look at together. 

\clearpage
\end{document}
